\section{Worked example: xl75\_multi}\label{sec:worked}

This section grounds the abstract machinery of \cref{sec:datamodel,sec:resolver,sec:diagnostics}
in a single concrete design from the test corpus,
\srcf{tests/data/designs/xl75\_multi.qrd}. The fixture is small---a nose cone, one body
tube, a fin set, and a coupler---yet it exercises all three seat kinds and, critically, it
is the design whose legacy \texttt{<offset>} encoding silently produced an interpenetrating
coupler. Tracing it end to end shows both how the station-pair model expresses each physical
relationship with the minimum of authored numbers and how the same model converts the
historical silent corruption into a hard, located failure. Every dimension quoted below is
read verbatim from the current fixture (it is still a version~\texttt{0.1} file carrying the
legacy CM-to-CM \texttt{<offset>} elements); the resolved geometry is the result the
station-pair resolver of \cref{sec:resolver} would produce after migration.

\subsection*{The root part and its frame}

The root is a solid \code{ConicalNoseCone} named \texttt{MultiNose} with $L = 0.30\,\text{m}$
and $\text{baseRadius} = 0.0395\,\text{m}$. In the settled longitudinal convention
($\zfwd$, see \cref{sec:philosophy}) the root is visited at the identity pose, so its aft
plane---the base rim---sits at world $z = 0$ and the tip is at world $z = +0.30$. Two
geometry-derived landmarks anchor everything downstream:

\begin{itemize}
  \item the aft station, $\{z = 0,\ r_{\text{outer}} = 0.0395\}$, and the tip station,
        $\{z = 0.30,\ r_{\text{outer}} \approx 0\}$;
  \item the closed-form outer skin, $\text{radiusOuterAt}(z) = 0.0395\,(1 - z/0.30)$, which
        the solid-host rule of \cref{sec:datamodel} treats as the host capacity because a
        solid cone has no bore ($\text{innerCapacityAt}(z) = \text{radiusOuterAt}(z)$).
\end{itemize}

The cone is also the design's single frame collision. Its centroidal tensor is computed in a
mid-length-origin, base-at-$+z$ frame by \src{model/parts/ConicalNoseCone.cpp}{49}, whose
\code{coneCmOffset} returns $z = L/2 - \bar h$ with $\bar h = L/4$ for a solid cone. That
internal $+z$ is the \emph{opposite} of the chosen forward direction. The resolver does not
flip the (correct, tested) tensor math; it applies the explicit forward-frame adapter of
\cref{sec:datamodel}, under which the solid cone's CM lies $\bar h$ forward of the base---that is,
$L/4 = 0.075\,\text{m}$ forward of the base (equivalently $3L/4$ aft of the tip), at world
$z = 0.075$. This single scalar, $\code{cmStationLocal} = 0.075$, is the only consumer of \src{model/parts/Part.h}{96}
(\code{getCenterMassOffset}) in the whole design, and it is the reason CM can fall out as a
derived quantity rather than being authored.

\tikzfig{worked-example}{The resolved \texttt{xl75\_multi} stack in the $\zfwd$ root frame.
The nose occupies $z\in[0,0.30]$; the body hangs aft over $z\in[-0.90,0]$; the fin set seats
on the body wall near the aft end; the coupler nests into the body fore bore and projects
$0.04\,$m forward of the bore mouth, where its outer radius exceeds the tapering nose skin.}{fig:worked}

\subsection*{The four relationships}

\paragraph{(1) Body abuts nose --- zero authored numbers.}
The body is a \code{BodyTube} (inner $0.0376$, outer $0.0395$, length $0.90$). It is attached
with the defaulted link, \seat{Abut}, $\text{parentStation01}=0$, $\text{childStation01}=1$,
$\text{gap}=0$: the child's fore plane meets the parent's aft plane. The resolver reads
$p = \text{nose}.\text{stationAt}(0) = \{z=0,\ r_{\text{outer}}=0.0395\}$ and
$c = \text{body}.\text{stationAt}(1) = \{z=0.90,\ r_{\text{outer}}=0.0395\}$, giving a child
aft-plane origin of $z = 0 + 0 - 0.90 = -0.90$. The body therefore spans world
$z \in [-0.90, 0]$, extending \emph{aft} ($-z$)---physically correct, the airframe hangs
below the nose. The Layer-1 seam check is clean: $|0.0395 - 0.0395| = 0$.

\paragraph{(2) Fin set seats on the body wall.}
The \code{FinSet} (root chord $0.10$, $\text{bodyRadius}=0.0395$, six fins) is attached
\seat{OnSurface} at $\text{parentStation01}=0.06$, $\text{childStation01}=0$,
$\text{gap}=0$. The resolver reads $p = \text{body}.\text{stationAt}(0.06)$, i.e. local
$z = 0.054$, world $z \approx -0.846$, with $r_{\text{outer}}=0.0395$, and seats the fin
root's aft end there. Layer~1 confirms radial compatibility: the fin's $\text{bodyRadius}$
of $0.0395$ equals the tube outer radius within tolerance. The true off-axis radial seat
is deferred to the 6-DOF work (\cref{sec:resolver,sec:sixdof}); because the fin set is
axisymmetric its mass sits on the longitudinal axis, so its 3-DOF placement is already exact.

\paragraph{(3) Coupler nests into the body fore bore.}
The coupler is a short \code{BodyTube} (inner $0.036$, outer $0.0376$, length $0.08$). Its
physical intent is ``insert its aft end $40\,$mm into the body's fore bore'', expressed as
\seat{NestInBore}, $\text{parentStation01}=1.0$, $\text{childStation01}=0.0$,
$\text{gap}=0.04$: the coupler's \emph{aft} plane is the inserted reference. The seat kind makes
the gap an insertion depth, so $\text{signedGap} = -0.04$ (the child moves aft, $-z$, into the
bore). The resolver reads $p = \text{body}.\text{stationAt}(1.0)$ at body-local $z = 0.90$ (the
fore bore mouth, world $z=0$) and $c = \text{coupler}.\text{stationAt}(0)$ at local $z = 0$,
giving a child aft-plane origin of $0.90 + (-0.04) - 0 = 0.86$ in the body frame, which composes
with the body's world origin $-0.90$ to world $z = -0.04$. The coupler thus spans world
$z \in [-0.04, +0.04]$: $0.04\,$m of it is correctly buried in the bore, but its fore $0.04\,$m
projects \emph{forward, past the body fore rim, into the nose region}, because the $0.08\,$m
coupler is longer than its $0.04\,$m insertion depth. Layer~1 passes---coupler outer radius
$0.0376 \le$ body inner radius $0.0376$, so it fits the bore radially---which is precisely why a
radial seam check alone is insufficient and the envelope sweep of \cref{sec:diagnostics} is
required.

\paragraph{(4) The forward poke-through, prevented.}
The coupler's forward $0.04\,$m occupies world $z \in [0, +0.04]$, beyond the body and inside
the nose's axial span. The Layer-2 envelope sweep resolves the host by the sorted-interval
query at each feature breakpoint:

\begin{itemize}
  \item Over $z \in [-0.04, 0]$ the host is the body bore, with capacity
        $r_{\text{inner}} = 0.0376$; the coupler's $0.0376 \le 0.0376$ holds, so this span is
        clean.
  \item Over $z \in [0, +0.04]$ the body is gone and the only covering envelope is the solid
        nose cone. The solid-host rule compares the coupler outer radius against the cone
        \emph{skin}: at the worst station $z = 0.04$,
        $\text{radiusOuterAt}(0.04) = 0.0395\,(1 - 0.04/0.30) = 0.0342$, while the coupler
        outer radius is $0.0376$. Since $0.0376 > 0.0342$, the coupler penetrates the cone
        skin by $0.0034\,$m.
\end{itemize}

\begin{hardfail}[Located poke-through: coupler through nose skin]
The sweep emits \cppinline|OverlapDiagnostic{offender = coupler, host = nose, zWorld = 0.04,|
\cppinline|penetration = 0.0034}| with the message ``coupler OD $0.0376$ exceeds nose skin
radius $0.0342$ at $z=0.04$ forward of the body rim (poke-through $0.0034\,$m)'', and sets
\cppinline|SolveResult.ok = false|. Per the gate in \cref{sec:diagnostics}, \emph{both}
consumers refuse the solve: \code{computeCompositeAt} throws/logs rather than integrate a
self-intersecting body, and \code{RocketMesh} renders the offender in the error colour and
surfaces the diagnostic rather than drawing a coupler buried in the cone. The failure is hard
and located---never silent. \fail
\end{hardfail}

\begin{table}[htbp]
\centering
\caption{The four \texttt{xl75\_multi} relationships, their station-pair links, and the
resolved result in the $\zfwd$ root frame.}\label{tab:worked}
\begin{tabularx}{\linewidth}{@{}l l X@{}}
\toprule
\textbf{Relationship} & \textbf{Link (seat; stations; gap)} & \textbf{Result} \\
\midrule
Body $\to$ nose &
\seat{Abut}; $0 \!\to\! 1$; gap $0$ &
Body origin $z=-0.90$; spans $[-0.90,0]$; seam $|0.0395-0.0395|=0$ \quad\ok \\
\addlinespace
Fins $\to$ body wall &
\seat{OnSurface}; $0.06 \!\to\! 0$; gap $0$ &
Root seats at world $z\approx-0.846$; $0.0395=0.0395$ \quad\ok \\
\addlinespace
Coupler $\to$ body bore &
\seat{NestInBore}; $1 \!\to\! 0$; gap $0.04$ &
Origin $z=-0.04$; spans $[-0.04,+0.04]$; OD $0.0376\le$ ID $0.0376$ \quad\ok (Layer~1) \\
\addlinespace
Coupler vs.\ nose skin &
(Layer-2 sweep over $[0,+0.04]$) &
At $z=0.04$, OD $0.0376 > $ skin $0.0342$; penetration $0.0034$ \quad\fail \\
\bottomrule
\end{tabularx}
\end{table}

\Cref{tab:worked} and \cref{fig:worked} together summarize the resolve: three relationships
that pass and one that is caught.

\subsection*{Why the legacy file failed silently}

In the version~\texttt{0.1} fixture the coupler is encoded as \cppinline|<offset z="0.49">|
relative to the body---a CM-to-CM displacement in the same $\zfwd$ frame. That offset slid
the coupler bodily \emph{forward} through the cone, and nothing checked it. Worse, the two
consumers disagreed on \emph{where} they drew it. The simulator accumulated the raw axial
offset, threading \cppinline|axialStation + pos.z()| through
\src{model/parts/Part.cpp}{224} in \code{accumulateAeroAt} and shifting child tensors to the
composite CM in \code{computeCompositeAt} (\src{model/parts/Part.cpp}{169}). The visualizer
walked the tree with \cppinline|station + offset| at \src{visualizer/RocketMesh.cpp}{472},
translating each \emph{geometric-centered} primitive (the destructure at
\src{visualizer/RocketMesh.cpp}{468}). Because a solid cone's geometric center and its CM
differ by exactly the $L/4$ offset of \src{model/parts/ConicalNoseCone.cpp}{49}, the two
renderings diverged by that cone offset \emph{on top of} the shared interpenetration error.
There was no signal at all: both pictures looked plausible, and a self-intersecting solid was
fed to the integrator.

\begin{keyinsight}
The station-pair model removes both faults at once. The stored intent is now ``nest $40\,$mm
into the body's fore bore'', and \seat{NestInBore} fixes the insertion \emph{direction}, so
the coupler can no longer be slid arbitrarily forward by an authored number. A single
resolver then produces one geometric placement that both consumers read, eliminating the
CM-versus-geometric-center divergence by construction. What remains of the original
mistake---the $0.04\,$m of coupler that still projects past the bore mouth---is caught by the
envelope sweep with the correct solid-host comparison and reported as a located failure.
Resolving the fixture cleanly is then a matter of physical intent: reduce the insertion gap
or shorten the coupler until its forward projection no longer exceeds the nose skin.
\end{keyinsight}
